\subsection{对角自同态}

定义 1. --- 设 \(\mathrm{B} = (e_i)_{i \in \mathrm{I}}\) 是 E 的一个标准正交基。称一个自同态 \(u\) 在基 B 中是对角的或相对于 B 是对角的，如果存在由 K 的元素组成的族 \(\lambda = (\lambda_i)_{i \in \mathrm{I}}\) ，使得 \(u(e_i) = \lambda_i e_i\) （对所有 \(i \in \mathrm{I}\) ）。

设 \(\mathrm{B} = (e_i)_{i \in \mathrm{I}}\) 是 E 的一个标准正交基。我们用 \(\mathcal{D}_{\mathrm{B}}(\mathrm{E})\) 表示 E 的相对于 B 为对角的自同态所成的集。这是 \(\mathcal{L}(\mathrm{E})\) 的一个闭的、幺的交换子代数。设 \(u \in \mathcal{D}_{\mathrm{B}}(\mathrm{E})\) 。族 \(\lambda = (\lambda_i)_{i \in \mathrm{I}}\) （满足 \(u(e_i) = \lambda_i e_i\) ）由 u 和基 B 唯一确定，称为 u 相对于 B 的特征值族。

假设 K = R，并把 E 与 \(\mathrm{E}_{(\mathbf{C})}\) 的一个子空间等同。设 \(\mathrm{B} = (e_i)_{i \in \mathrm{I}}\) 是 E 的一个标准正交基。用 \(\mathrm{B}_{(\mathbf{C})}\) 表示标准正交基 \((1 \otimes e_i)_{i \in \mathrm{I}}\) （\(\mathrm{E}_{(\mathbf{C})}\) 的）（参见 EVT, V, p. 29, 推论 1）。映射 \(u \mapsto u_{(\mathbf{C})}\) 诱导从 \(\mathcal{D}_{\mathrm{B}}(\mathrm{E})\) 到 \(\mathcal{D}_{\mathrm{B}_{(\mathbf{C})}}(\mathrm{E}_{(\mathbf{C})})\) 中的一个单的代数态射；它的像是 \(u \in \mathcal{D}_{\mathrm{B}_{(\mathbf{C})}}(\mathrm{E}_{(\mathbf{C})})\) （满足 \(u(\mathrm{E}) \subset \mathrm{E}\) ）所成的集；换言之，即基 \(\mathrm{B}_{(\mathbf{C})}\) 中特征值为实的对角自同态 u 所成的集。

对每个 \(i \in \mathrm{I}\) ，用 \(p_i\) 表示 E 的以 \(\mathrm{Ke}_i\) 为像的正交投影。我们有 \(\| p_i \| = 1\) （对所有 \(i \in \mathrm{I}\) ）。自同态 \(p_i\) 在基 B 中是对角的，其特征值族是族 \((\delta_{ij})_{j\in\mathrm{I}}\) （克罗内克符号，参见 A, II, p. 24）。

设 \(\mathrm{B}^{\prime}=(f_{i})_{i\in\mathrm{I}}\) 是 E 的一个标准正交基，\(u:E\to E\) 是满足 \(u(f_{i})=e_{i}\) 的等距同构。则 \(\mathcal{D}_{\mathrm{B}^{\prime}}(\mathrm{E})=u^{-1}\mathcal{D}_{\mathrm{B}}(\mathrm{E})u\) 。

我们给 I 赋以离散拓扑，并用 \(\mathcal{B}(\mathrm{I}) = \mathcal{C}_{b}(\mathrm{I}; \mathrm{K})\) 表示 I 上取值于 K 的有界函数所成的幺巴拿赫代数（I, p. 17, 例 2）；若 K = C，它是一个星代数（I, p. 102, 例 2）。

引理 1. --- 设 \(u \in \mathcal{D}_{\mathrm{B}}(\mathrm{E})\) ，\(\lambda = (\lambda_i)_{i \in \mathrm{I}}\) 是它的特征值族。族 \(\lambda\) 有界，且有 \(\sup_i |\lambda_i| = \| u\|\) 。

有 \(|\lambda_i| \leqslant \|u\|\) （对所有 \(i\) ），从而 \(\sup_{i \in I} |\lambda_i| \leqslant \|u\|\) 。此外，由于

\[
\| u (x) \| ^ {2} = \sum_ {i \in \mathrm{I}} | \lambda_ {i} | ^ {2} | \langle e _ {i} | x \rangle | ^ {2} \leqslant \left(\sup _ {i \in \mathrm{I}} | \lambda_ {i} | ^ {2}\right) \| x \| ^ {2}
\]

对所有 \(x \in \mathrm{E}\) 成立（EVT, V, p. 22, 命题 5），由此得出 \(\| u\| = \sup |\lambda_i|\) 。

命题 1. --- a) 映射 \(\alpha\) （从 \(\mathcal{D}_{\mathrm{B}}(\mathrm{E})\) 到 \(\mathcal{B}(\mathrm{I})\) 中的）把对角自同态 \(u\) 对应到 \(u\) 的特征值族，它是 K 上的巴拿赫代数的一个等距同构。若 K = C，它是对合代数的一个态射。

\begin{enumerate}
\def\labelenumi{\alph{enumi})}
\setcounter{enumi}{1}
\item
  对任何有界族 \(\lambda = (\lambda_i)_{i\in \mathrm{I}}\) ，族 \((\lambda_ip_i)_{i\in \mathrm{I}}\) 在赋以逐点收敛拓扑的空间 \(\mathcal{L}(\mathrm{E})\) 中是可和的，且该族之和是唯一的映射 \(u\in \mathcal{D}_{\mathrm{B}}(\mathrm{E})\) ，它满足 \(\alpha (u) = \lambda\) ；
\item
  设 \(u \in \mathcal{D}_{\mathrm{B}}(\mathrm{E})\) ，\(\lambda\) 是它相对于 B 的特征值族。\(u\) 的谱是 \(\lambda\) 的值集在 K 中的闭包；
\item
  若 K = C，则从 \(\mathcal{C}(\mathrm{Sp}(u))\) 到 \(\mathcal{L}(\mathrm{E})\) 中的连续函数演算映射把连续函数 \(f \in \mathcal{C}(\mathrm{Sp}(u))\) 对应到自同态 \(\alpha(f \circ \lambda)\) （\(\mathcal{D}_{\mathrm{B}}(\mathrm{E})\) 中的）。
\end{enumerate}

由引理 1，在 B 中为对角的自同态的特征值族是有界的，且如此定义的从 \(\mathcal{D}_{\mathrm{B}}(\mathrm{E})\) 到 \(\mathcal{B}(\mathrm{I})\) 中的映射是幺巴拿赫代数的一个连续等距态射。

设 \(i \in I\) 。对每个 \(j \in I\) ，我们有 \(\langle u^{*}(e_{i}) | e_{j} \rangle = \lambda_{j} \langle e_{i} | e_{j} \rangle = \langle \overline{\lambda_{j}} e_{i} | e_{j} \rangle\) 。由此得出 \(u^{*}(e_{i}) = \overline{\lambda}_{i} e_{i}\) 。于是 u 的伴随是基 B 中的对角自同态，其特征值族是 \((\overline{\lambda}_{i})_{i \in I}\) 。断言 a) 由此得出。

设 \(\lambda = (\lambda_i)_{i\in \mathrm{I}}\in \mathcal{B}(\mathrm{I})\) 。对每个 \(x\in \mathrm{E}\) ，族 \((|\langle e_i|x\rangle |^2)_{i\in \mathrm{I}}\) 可和，其和为 \(\| x\|^2\) （EVT, V, p. 22, 命题 5），从而，对 I 的每个有限子集 J：

\[
\begin{array}{r l r} & & {\left\| \sum_ {i \in \mathrm{J}} \lambda_ {i} p _ {i} (x) \right\| ^ {2} = \sum_ {i \in \mathrm{J}} | \lambda_ {i} | ^ {2} | \langle e _ {i} | x \rangle | ^ {2} \leqslant \Bigl (\underset {i \in \mathrm{I}} {\sup} | \lambda_ {i} | ^ {2} \Bigr) \sum_ {i \in \mathrm{J}} | \langle e _ {i} | x \rangle | ^ {2}} \\ & & {\leqslant \Bigl (\underset {i \in \mathrm{I}} {\sup} | \lambda_ {i} | \Bigr) ^ {2} \| x \| ^ {2}.} \end{array}
\]

因此，族 \((\lambda_{i}p_{i})\) 在赋以简单收敛拓扑的 \(\mathcal{L}(\mathrm{E})\) 中可和。它的和 \(u_{\lambda}\) 满足 \(u_{\lambda}(e_{i}) = \lambda_{i}e_{i}\) ；于是它是基 B 中的对角自同态，以 \(\lambda\) 为特征值族。断言 b) 由此得出。

最后的断言由 \(a\) )、I, p. 17 的例 3 与 I, p. 111 的例 4 得出。

注记. --- 巴拿赫代数 \(\mathcal{D}_{\mathrm{B}}(\mathrm{E})\) 是 \(\mathcal{L}(\mathrm{E})\) 的一个极大交换闭子代数。事实上，设 u 是与 \(\mathcal{D}_{\mathrm{B}}(\mathrm{E})\) 交换的 E 的自同态。设 \(i \in I\) 。由于正交投影 \(p_{i}\) 在基 B 中是对角的，我们有 \(p_{i}(u(e_{i})) = u(p_{i}(e_{i})) = u(e_{i})\) ，这蕴含 \(u(e_{i})\) 与 \(e_{i}\) 成比例。于是 u 在基 B 中是对角的。

若 E 是无穷维的，则存在 \(\mathcal{L}(\mathrm{E})\) 的极大交换对合子代数，它们不同构于 \(\mathcal{D}_{\mathrm{B}}(\mathrm{E})\) （IV, p. 314, 习题 5）。

命题 2. --- 设 \(u \in \mathcal{D}_{\mathrm{B}}(\mathrm{E})\) ，\(\lambda = (\lambda_i)_{i \in \mathrm{I}}\) 是它的特征值族。

\begin{enumerate}
\def\labelenumi{\alph{enumi})}
\tightlist
\item
  下列条件等价：
\end{enumerate}

\begin{enumerate}
\def\labelenumi{(\roman{enumi})}
\item
  有 \(\lambda \in \mathcal{C}_{0}(\mathrm{I};\mathrm{K})\) ；
\item
  自同态 u 是紧的；
\item
  族 \((\lambda_{i}p_{i})_{i\in\mathrm{I}}\) 在巴拿赫空间 \(\mathcal{L}(\mathrm{E})\) 中可和。其和这时等于 u。
\end{enumerate}

\begin{enumerate}
\def\labelenumi{\alph{enumi})}
\setcounter{enumi}{1}
\tightlist
\item
  假设 \(u\) 是紧的，用 \(\Lambda\) 表示 \(\lambda\) 的值集。由 \(i \in I\) （使 \(\lambda_i \neq 0\) ）组成的集是可数的。若 E 是无穷维的，我们有 \(\mathrm{Sp}_{\mathrm{s}}(u) = \Lambda - \{0\}\) 与 \(\mathrm{Sp}(u) = \Lambda \cup \{0\}\) 。若 E 是有限维的，则 \(\mathrm{Sp}_{\mathrm{s}}(u) = \mathrm{Sp}(u) = \Lambda\) 。
\end{enumerate}

由引理 1，对 I 的每个有限子集 J，我们有 \(\| \sum_{j\in J}\lambda_jp_j\| = \sup_{j\in J}|\lambda_j|\) 。由此，按柯西判别准则，族 \((\lambda_ip_i)\)

在 \(\mathcal{L}(\mathrm{E})\) 中可和当且仅当族 \(\lambda\) 在无穷远处趋于 0，这蕴含条件 (i) 与 (iii) 等价。

条件 (iii) 蕴含 \(u\) 是紧的（III, p. 4, 命题 2 的推论）。反之，假设自同态 \(u\) 是紧的。由于族 B 在 E 中有界，它在 \(u\) 下的像在 E 中相对紧（III, p. 2），从而是准紧的。设 \(\varepsilon > 0\) ，设 \(J\) 是 I 的一个有限子集，使得 B 在 \(u\) 下的像包含于半径为 \(\varepsilon\) 、心为 \(u(e_j)\) （\(j \in \mathrm{J}\) ）的球之并。设 \(i \in \mathrm{I} - \mathrm{J}\) 。存在 \(j \in \mathrm{J}\) 使 \(\| u(e_i) - u(e_j) \| \leqslant \varepsilon\) ，从而

\[
\left| \lambda_ {i} \right| ^ {2} \leqslant \left| \lambda_ {i} \right| ^ {2} + \left| \lambda_ {j} \right| ^ {2} = \left\| u \left(e _ {i}\right) - u \left(e _ {j}\right) \right\| ^ {2} \leqslant \varepsilon^ {2}.
\]

因此，我们有 \(\lambda \in \mathcal{C}_0(\mathrm{I};\mathrm{K})\) ，即条件 (i)。

最后，u 的谱及其敏感谱利用 \(\lambda\) 由命题 1, c) 与 III, p. 90 命题 5 计算。

注记. --- 当 I 无限时，条件 \(\lambda\in\mathcal{C}_{0}(\mathrm{I};\mathrm{K})\) 也可表述为“族 \(\lambda\) 沿 I 的有限子集的余集的滤子趋于 0”。

